\title{Byte Latent Transformer: Patches Scale Better Than Tokens}

\begin{abstract}
We present the Byte Latent Transformer ({\textsc BLT}), a novel architecture for byte-level large language models (LLMs) that, for the first time, achieves performance comparable to models relying on tokenization, while offering notable gains in both inference speed and model robustness. {\textsc BLT} processes input bytes by grouping them into variable-length patches, which constitute the fundamental units for computation. These patches are delineated according to the entropy associated with predicting the subsequent byte, allowing the model to allocate more resources and capacity to sections of higher informational complexity. We conduct the inaugural flop-controlled scaling investigation of byte-level architectures, scaling up to 8B parameters and 4T training bytes. Our findings establish that scaling LLMs directly on raw byte streams—without dependence on a pre-defined vocabulary—is both practical and effective. Training and inference benefit from the adaptive selection of longer patches for predictable data segments, contributing to improved efficiency as well as qualitative advances in reasoning capabilities and generalization to less frequent patterns. In summary, for equivalent inference budgets, {\textsc BLT} achieves superior scaling properties compared to tokenization-based approaches, leveraging concurrent increases in both patch lengths and model size.
\end{abstract}

\section{Introduction}
\label{section:intro}

We present the Byte Latent Transformer~(\textbf{BLT}), an architecture that operates directly on raw byte sequences, completely bypassing the need for tokenization. For the first time, BLT achieves parity with tokenization-based models at scale, while also offering notable gains in both efficiency and robustness (\S\ref{sec:robustness}). In contrast, standard large language models (llms) are predominantly trained in an end-to-end manner with the exception of the tokenization stage—an ad hoc preprocessing step that converts bytes into a predefined vocabulary of tokens. This reliance on tokenization introduces compression biases into the representation of text, resulting in limitations such as sensitivity to domain and modality~\citep{dagangetting}, heightened vulnerability to input perturbations~(\S\ref{sec:robustness}), insufficient encoding of orthographic features~\citep{edman2024cute}, and systemic disadvantages for some languages in multilingual settings~\citep{liang2023xlm,petrov2024language,limisiewicz2024myte}.

Tokenization has traditionally played a vital role, as training large language models (LLMs) directly on raw bytes incurs excessive computational costs at scale, largely due to the increased sequence lengths~\citep{xue2022byt5}. Some earlier studies have addressed this challenge by introducing more efficient variants of self-attention~\citep{el2020characterbert, clark2022canine} or leveraging architectures that eliminate attention mechanisms altogether~\citep{wang2024mambabyte}~(\S\ref{section:related_work}). Nevertheless, these solutions are predominantly effective for training \textit{small models}. When scaling up, the overall computation in a Transformer is primarily dictated by the expansive feed-forward network layers that process every single byte, rather than by the attention component itself.

In order to optimize compute allocation, we introduce an adaptive and learnable technique for clustering bytes into \textit{patches}~(\S\ref{section:patching}), alongside a novel model architecture that integrates both byte-level and patch-level representations. Distinct from tokenization, {\textsc BLT} does not rely on a predetermined set of patch vocabulary. Instead, it utilizes efficient trainable encoder and decoder components to transform arbitrary byte groupings into latent patch embeddings. Our results demonstrate that this approach allows for a \textit{more} effective use of computational resources compared to traditional tokenization-based frameworks.

Large language models that rely on tokenization assign an equal amount of computational resources to every token. This approach prioritizes simplicity at the expense of optimal performance, as the heuristics used to generate tokens may not align with the true complexity involved in making predictions. The core principle behind our architecture is that computation should be distributed adaptively, targeting regions where it is most necessary. For instance, utilizing a large transformer network is generally unnecessary for tasks like predicting word endings, which are typically straightforward and present low entropy compared to predicting the initial word in a new sentence. In {\textsc BLT}'s architecture~(\S\ref{section:architecture}), this idea is operationalized through the deployment of three transformer modules: two lightweight byte-level \textit{local models} complemented by a robust global \textit{latent transformer}~(\autoref{fig:arch_high_level}).
{\textsc BLT} addresses the challenge of grouping bytes into patches—and thus adaptively allocating compute—by segmenting sequences according to the entropy of upcoming byte predictions. This results in groupings of bytes with comparable information density, supporting efficient, context-sensitive processing.

We introduce the first scaling analysis of byte-level models that is controlled by floating point operation (flop) counts, exploring models with up to 8B parameters and trained on 4T bytes of data. Our findings demonstrate that large-scale, end-to-end training from raw bytes, without relying on fixed-vocabulary tokenization, is feasible. Notably, {\textsc BLT} achieves parity in training performance, as measured by controlled flop usage,\footnote{Our assessment of computational demand is based on the total number of Floating Point OPerations (flops) performed.} with Llama 3, while attaining up to 50\% reduction in inference flops (\S\ref{section:scaling}). Additionally, we establish that operating directly on raw bytes yields notable gains in representing rare and infrequent data distributions. {\textsc BLT} models demonstrate greater resilience to input noise compared to tokenization-dependent models and exhibit superior abilities at the character level, as seen in evaluations involving orthographic processing, phonological analysis, and low-resource machine translation (\S\ref{sec:robustness}). Furthermore, the {\textsc BLT} framework allows simultaneous scaling of both model and patch sizes without exceeding a fixed inference flop allocation. Increased patch lengths, on average, lead to computational savings that can be redirected to expand the global latent transformer since it executes less frequently. Through inference-flop constrained scaling studies (\autoref{fig:fixed_inference_scaling}), we find that this architecture delivers markedly improved scaling dynamics compared to those based on tokenization.

To conclude, this work offers several key contributions: 1) We present {\textsc BLT}, a byte latent llm framework that adaptively distributes computation to enhance flop efficiency; 2) We provide evidence that {\textsc BLT} attains parity with Llama 3 in terms of training flop usage at scales up to 8B parameters, with the flexibility to achieve up to 50\% gains in flop efficiency by allowing minimal drops in evaluation performance; 3) {\textsc BLT} introduces a novel approach for scaling llms, enabling growth in model size without exceeding a predetermined inference compute budget; 4) Our results illustrate that {\textsc BLT} models exhibit heightened resilience to noisy inputs and a finer sensitivity to sub-word features in the input, which token-centric llms tend to overlook. The codebase for both training and inference with {\textsc BLT} is publicly accessible at~\url{https://github.com/facebookresearch/blt}.

\section{Patching: From Individual Bytes to Groups of Bytes}
\label{section:patching}

Dividing a byte sequence into \textit{patches} enables {\textsc BLT} to adaptively assign computational resources according to contextual requirements. Figure~\ref{fig:patching_types} illustrates various strategies for partitioning bytes into patches. More precisely, a patching function $f_p$ takes a sequence of bytes $\pmb{x}=\{x_i,|i=1,\ldots n\}$ of length $n$ and produces a shorter sequence of $m < n$ patches, denoted $\pmb{p}=\{p_j|j=1,\ldots,m\}$, by assigning each $x_i$ to a value in \{0,1\}; a value of 1 signifies the start of a new patch. Whether considering token-based or patch-based approaches, the principal computational burden in data processing stems from the total number of steps the core Transformer must perform. In the context of {\textsc BLT}, this equates to the total number of patches created by the selected patching function. Therefore, the \textit{patch size}—the typical number of bytes within a single patch—serves as the primary determinant of computational expense during both training and inference with a chosen patching scheme~(\S\ref{section:flops}).

The subsequent sections present three specific patching functions: partitioning data into patches of a constant byte count~(\S\ref{section:static-patch}), using whitespace as patch boundaries~(\S\ref{section:white-patch}), and leveraging byte-level language model entropy for adaptive patch segmentation~(\S\ref{section:dyn-patch}). The discussion concludes by covering incremental patching techniques and contrasting patching with conventional tokenization approaches~(\S\ref{section:bpe-patch}).

\subsection{Strided Patching Every K Bytes}
\label{section:static-patch}

One of the simplest methods for grouping bytes involves dividing them into uniform patches of size $k$, as utilized in MegaByte~\citep{yu2023megabyte}. This fixed-stride approach is straightforward to apply during both training and inference, allows for easy adjustments to the average patch size, and thus offers a direct way to manage computational cost. Despite these benefits, there are notable limitations to this patching strategy. Firstly, computation is not allocated adaptively to different parts of the sequence: valuable processing steps may be wasted when predicting uninformative segments such as whitespace in code, while highly informative bytes, such as those representing mathematical content, may not receive adequate attention. Secondly, this method introduces inconsistent and non-contextual patch segmentation for identical byte sequences, leading to issues such as the inconsistent splitting of identical words.

\subsection{Space Patching}
\label{section:white-patch}

According to \citet{slagle2024spacebyte}, an effective modification to strided patching is introduced whereby new patches are initiated after encountering any space-like byte\footnote{
    Space-like bytes refer to any byte not corresponding to a latin letter, digit, or utf-8 continuation byte. Furthermore, each patch is required to include at least one byte that is not space-like.}, which tend to align with the natural boundaries of linguistic units across many languages. The approach, termed Space patching, dedicates a latent transformer computation (thus increasing flops) to each word, enabling consistent patching of words throughout different sequences and assigning greater computational resources to positions where challenging predictions arise—typically immediately after spaces. As an illustration, generating the initial byte of the answer in the query ``Who composed the Magic Flute?\uline{\hspace{.6em}}'' involves a considerably higher degree of difficulty compared to producing subsequent bytes following ``M'', since this first character constrains the prediction space, thereby rendering the remainder (i.e., ``Mozart'') easier to predict. Nevertheless, space patching faces difficulties in adapting to all language contexts and application domains, with a significant limitation being its inability to dynamically adjust patch sizes. In the following, we present an alternative patching strategy based on the observation that the first bytes of words are generally the most unpredictable, while simultaneously providing an intuitive way to manage patch dimensions.

\subsection{Entropy Patching: Using Next-Byte Entropies from a Small Byte LM}
\label{section:dyn-patch}

Instead of utilizing a rule-based heuristic like whitespace, our method employs a data-driven strategy to detect locations where next-byte predictions are most uncertain. We propose \textit{entropy patching}, which determines patch boundaries based on entropy measures.

Specifically, we fit a compact auto-regressive language model at the byte level using the training set from {\textsc BLT}. This model is then used to calculate next byte entropies according to the language model probability distribution $p_{e}$ defined over the byte vocabulary $\mathcal{V}$:

\begin{equation}
H(x_i) = \sum_{v \in \mathcal{V}} p_{e}(x_i=v|\pmb{x}_{<i}) \log p_{e}(x_i=v|\pmb{x}_{<i})
\end{equation}

We utilize two distinct techniques to determine patch boundaries using entropies $H(x_i)$. The initial method selects locations where the entropy surpasses a predefined global threshold, as depicted in~\autoref{fig:patching}. The alternative strategy detects locations where the entropy is elevated relative to the preceding value. This latter method can be viewed as pinpointing locations where the entropy deviates from an otherwise roughly monotonic decrease within a patch.

\begin{align*}
    \text{Global Constraint}\hspace{1cm}H(x_t)& > \theta_g\\
    \text{Approx. Monotonic Constraint}\hspace{1cm}H(x_t)& - H(x_{t-1}) > \theta_r
\end{align*}

Patch boundaries are determined through a simple preprocessing procedure that occurs as part of dataloading. This approach contrasts with~\citet{nawrot-etal-2023-efficient}, where a classifier is utilized to predict patch boundaries based on entropy estimates. In our experiments~(\S\ref{section:experiments}), we evaluate and compare both techniques for differentiating low- and high-entropy bytes.

\subsection{The Byte-Pair Encoding (BPE) Tokenizer and Incremental Patching}
\label{section:incremental-patching}
\label{section:bpe-patch}

A large number of contemporary llms, such as our baseline Llama 3, employ a subword tokenization approach like bpe~\citep{gage1994new, sennrich-etal-2016-neural}. Throughout this work, we use the term ``tokens'' to denote byte groupings selected from a \textit{finite} vocabulary established before training begins. In contrast, ``patches'' denote dynamically aggregated sequences that do not rely on a fixed vocabulary. One key distinction between tokens and patches is that tokens prevent the model from directly accessing the original byte-level features.

A significant advancement introduced by {\textsc BLT} compared to models relying on tokenization lies in its reconfiguration of the relationship between vocabulary size and computational requirements. Conventional large language models (LLMs) encounter a constraint: expanding the vocabulary generally results in larger average token sizes, thereby reducing the number of processing steps. However, this also increases the dimensionality of the model's output layer, leading to heightened computational and memory costs. This dynamic imposes restrictions on how much tokenization-based methods can flexibly adjust token granularity and inference efficiency. For instance, in moving from Llama 2 to Llama 3, the average token length grows from 3.7 to 4.4 bytes, but this improvement necessitates expanding the embedding table’s size by a factor of four.

During the generation process, {\textsc BLT} must determine at each position within the byte sequence whether it is at a patch boundary, since this affects if additional computation via the Latent Transformer is required. Importantly, this determination must be made without relying on knowledge of the bytes yet to be generated; segmentation into patches must be conducted solely based on already generated bytes. In formal terms, a patching function $f_p$ is said to have the property of incremental patching if the following holds:
$$f_p(\pmb{x}_{<i}) = f_p(\pmb{x})_{<i}$$
The bpe method does not possess the property of incremental patching, as the tokenization of any given prefix can vary depending on subsequent bytes in the sequence. Consequently, it does not fulfill the condition outlined above\footnote{Introducing a dedicated delimiter token to mark patch boundaries can enable bpe to function incrementally, but this comes at the cost of a longer byte sequence.}.

\section{{\textsc BLT} Architecture}
\label{section:architecture}
The {\textsc BLT} framework consists of a large, global autoregressive language model that operates on patch-level representations. Supplementing this are two more compact, local models: one that encodes byte sequences into patch representations and another that reconstructs bytes from these patches~(\autoref{fig:arch_high_level}).

\subsection{Latent Global Transformer Model}

\textit{The Latent Global Transformer} refers to an autoregressive transformer model $\mathcal{G}$, composed of $l_{\mathcal{G}}$ layers, which transforms an input sequence of latent patch representations $p_j$ into a corresponding sequence of output patch representations $o_j$. In this work, we assign the subscript $j$ to indicate individual patches, while $i$ denotes bytes. The global model employs a block-causal attention mask~\citep{dubey2024llama}, ensuring that attention is confined to patches up to and including the current one within the same document. Notably, this component is responsible for the majority of computational cost (flops) both during pre-training and inference. Consequently, selectively activating this module offers a mechanism to manage and adapt computational expenditure across different sections of the input and output, depending on their respective complexities.

\subsection{Local Encoder}
\label{section:local}

\textit{The Local Encoder Model}, represented as $\mathcal{E}$, employs a transformer architecture with a substantially smaller number of layers ($l_{\mathcal{E}} << l_{\mathcal{G}}$), focusing on rapid and effective conversion of input byte sequences $b_i$ into rich patch-level representations $p_j$. Unlike standard transformers, this model integrates a cross-attention mechanism subsequent to each transformer block, which serves to aggregate byte-level information into patch representations~(\autoref{fig:crossattn}). 

Initially, byte inputs $b_i$ are projected into vector space via an embedding matrix of shape $\mathbb{R}^{256 \times h_{\mathcal{E}}}$, producing embeddings $x_i$. Optionally, these representations may be further enriched through the incorporation of hash-embeddings as described in~(\S\ref{sec:hash-emb}). The embedded tokens are then passed through an interleaved sequence of transformer layers and cross-attention modules~(\S\ref{sec:enc-cross-attn}), progressively yielding patch representations $p_i$ that are subsequently input to the global transformer $\mathcal{G}$. 

Within each transformer block, a \textit{local block causal} attention mask is applied: every byte token has access only to a sliding window of $w_{\mathcal{E}}$ previous bytes. This windowing can overlap multiple patch boundaries, though it does not extend across document limits. Subsequent sections discuss the construction of the embedding space and elaborate on the workings of the cross-attention mechanism.

\subsubsection{Encoder Hash n-gram Embeddings}
\label{sec:ngram-emb}
\label{sec:hash-emb}
An essential aspect of building effective and descriptive representations at every step $i$ involves encoding contextual information from the preceding bytes. In {\textsc BLT}, this is addressed by representing not only the single byte $b_i$ at position $i$, but also by embedding it as a component within a broader byte n-gram context. For each step $i$, we begin by generating byte-grams

\begin{align}
    g_{i,n}=\{b_{i-n + 1},\ldots, b_i\}
\end{align}

for every byte position $i$ and for $n$ ranging from three to eight.\footnote{
Byte-grams of length $n$ or greater are not considered if $i<n$. }

Subsequently, we propose the use of hash $n$-gram embeddings, which utilize a hash function to assign each byte $n$-gram to an entry in an embedding table $E_{n}^{hash}$ of constant size, separately for each $n \in\{3,4,5,6,7,8\}$~\citep{bai2010learning}. The resulting vector is combined with the original byte embedding, then normalized prior to serving as input to the local encoder component. Formally, the enriched embedding is computed as follows:

\begin{align}
e_i &= x_i + \sum_{n = 3,...,8}E_{n}^{hash}(\text{Hash}(g_{i,n})) \\
\text{where, Hash}(g_{i,n}) &= \text{RollPolyHash}(g_{i,n}) \% |E_{n}^{hash}|
\end{align}

We adjust $e_i$ by dividing it by the total number of $n$-gram sizes incremented by one, employing RollPolyHash as specified in Appendix~\ref{appendix:rollpolyhash}. Section~\ref{section:ablations} examines the impact of varying both $n$ and the embedding table size on flop-controlled scaling laws by ablating the $n$-gram hash embeddings. Beyond hash-based $n$-gram embeddings, we also evaluated frequency-based $n$-gram embeddings; further information on these experiments can be found in Appendix~\ref{appendix:ngrams}.

\subsubsection{Encoder Multi-Headed Cross-Attention}
\label{sec:enc-cross-attn}
Our approach largely mirrors the input cross-attention mechanism utilized in the Perceiver architecture~\citep{jaegle2021perceiver}, with the principal distinction that our latent variables are dynamic, each associated with representations of variable patches rather than relying on a predetermined set of latent representations~(\autoref{fig:crossattn}). Each latent variable restricts its attention exclusively to the bytes that constitute its associated patch. Within this module, every patch $p_j$ is assigned a query vector, initially formed by aggregating the byte-level features linked to patch $p_j$. This is subsequently mapped through a linear transformation, denoted as $\mathcal{E}_{C} \in \mathbb{R}^{h_{\mathcal{E}} \times (h_{\mathcal{E}} \times U_{\mathcal{E}})}$, where $U_{\mathcal{E}}$ designates the number of encoder cross-attention heads. Precisely, for a patch $p_j$ with byte sequence $f_{\text{bytes}}(p_j)$, the calculation proceeds as follows:

\begin{align}
P_{0,j} &= \mathcal{E}_{C}(f_\text{bytes}((p_j)), f~ \text{is a pooling function} \\
P_l &= P_{l-1} + W_o\left(\text{softmax}\left(\frac{QK^T}{\sqrt{d_k}}\right)V\right) \\
\text{where } Q_j &= W_q(P_{l-1,j}), K_i = W_k(h_{l-1,i}), V_i = W_v(h_{l-1,i}) \\
h_l &= \text{Encoder-Transformer-Layer}_l(h_{l-1})
\end{align}

Here, $P \in \mathbb{R}^{n_p \times h_{\mathcal{G}}}$ denotes the set of $n_p$ patch representations designated for processing by the global model. This set is initialized by aggregating the byte embeddings $e_i$ that correspond to each patch $p_j$. The matrices $W_q$, $W_k$, $W_v$, and $W_o$ serve as the projection operators for the queries, keys, values, and outputs, respectively. For the keys and values, these projections are computed from the byte representations $h_i$ from the preceding layer (which are $e_i$ in the first layer). Our approach incorporates a patch-specific masking scheme, such that each query $Q_j$ can only attend to the keys and values associated with the bytes within patch $j$. Since multi-headed attention is applied over $Q, K$, and $V$, and patch representations typically have a higher dimensionality ($h_{\mathcal{G}}$) than $h_{\mathcal{E}}$, we represent $P_l$ as several attention heads of dimension $h_{\mathcal{E}}$ for the cross-attention step, and subsequently concatenate these head outputs to recover the $h_{\mathcal{G}}$ dimensions. Furthermore, pre-LayerNorm is applied to the queries, keys, and values, and no positional encodings are included in this cross-attention framework. The module is concluded with a residual connection that wraps around the cross-attention block.

\subsection{Local Decoder} 

The local decoder $\mathcal{D}$, mirroring the architecture of the local encoder, is also a compact transformer-based network comprising $l_{\mathcal{D}}$ layers, where $l_{\mathcal{D}} << l_{\mathcal{G}}$. Its role is to convert a series of global patch representations $o_j$ into the original byte sequence $y_i$. By leveraging the hidden states generated by the local encoder from the byte sequence, the local decoder predicts each byte conditioned on those previously generated. This process involves stacking $l_{\mathcal{D}}$ layers that alternate between cross-attention and standard transformer operations. Notably, the cross-attention module precedes the transformer block within each layer, facilitating the transformation from patch representations to byte-level representations, after which the transformer processes this byte sequence.

\subsubsection{Decoder Multi-headed Cross-Attention} 

Within the decoder’s cross-attention mechanism, the assignment of queries and key/value pairs is reversed: byte-representations function as the queries, while patch representations serve as the keys and values. At the outset, byte-representations for cross-attention are seeded using the byte embeddings produced by the final encoder layer, denoted as $h_{l_{\mathcal{E}}}$. For each following layer $l$, the byte-representations $d_{l,i}$ are determined as follows:

\begin{align}
D_0 &= h_{l_{\mathcal{E}}} \\
B_l &= D_{l-1} + W_o\left(\text{softmax}\left(\frac{QK^T}{\sqrt{d_k}}\right)V\right), \\
  &\text{where } Q_i = W_q(d_{l-1,i}), K_i = W_k(\mathcal{D}_C(o_j)), V_i = W_v(\mathcal{D}_C(o_j)) \\
  D_l &= \text{Decoder-Transformer-layer}_l(B_l)
\end{align}

In this context, $W_k$ and $W_v$ serve as the projection matrices for keys and values, acting via a linear transformation and split operation $\mathcal{D}_C$ applied to the final patch representations $o_j$ generated by the global model. The query projection matrix $W_q$ is applied to the byte-level representations $d_{l-1}$ from the preceding decoder transformer layer (or $h_{l_{\mathcal{E}}}$ in the case of the initial layer). The output projection is performed by $W_o$. Consequently, $B$ has shape $\mathbb{R}^{h_{\mathcal{D}} \times n_b}$, where $n_b$ denotes the total output bytes. The next-stage decoder representations, $D_l$, are then calculated by feeding the output $B$ from the cross-attention block into a decoder transformer layer. Following the methodology used in local encoder cross-attention, our design incorporates multi-head attention, applies pre LayerNorms, excludes positional embeddings, and includes a residual pathway around the cross-attention module.

\section{Experimental Setup}
\label{section:experiments}
We implemented a series of rigorously controlled experiments to benchmark {\textsc BLT} against models employing standard tokenization, ensuring that sequence context length does not unintentionally favor {\textsc BLT}.

\subsection{Pre-training Datasets} In this study, all tested model configurations are initially trained on two main datasets: (1) The Llama 2 dataset~\citep{touvron2023llama}, containing 2 trillion tokens sourced from an assortment of publicly available corpora, which undergo thorough cleaning and filtering to enhance data integrity; and (2) {\textsc BLT}-1T, a novel collection of 1 trillion tokens also compiled from multiple open sources, augmented by a selection of pre-training material shared by Datacomp-LM~\citep{li2024datacomp}. The Llama 2 dataset supports scaling law analyses for optimizing token counts as outlined in~\cite{dubey2024llama}, guiding the choice of architecture for {\textsc BLT}. In contrast, {\textsc BLT}-1T serves as the basis for full pre-training, facilitating downstream evaluation against Llama 3. Neither dataset contains any data obtained from Meta platforms or services. Additionally, for establishing baselines in tokenizer-centric models, we utilize the Llama 3 tokenizer with a 128K token vocabulary, which demonstrated superior baseline outcomes compared to the Llama 2 tokenizer throughout our tests.

\subsection{Entropy Model} For our experiments, the entropy model is implemented as a byte-level language model, utilizing the identical training distribution employed for the full {\textsc BLT} model. Unless stated otherwise, our default architecture comprises a transformer with 100M parameters, 14 layers, and a hidden size of 512, featuring sliding window attention spanning 512 bytes. All other hyperparameters align with those used for our local and global transformer models. Various model sizes, receptive field lengths, and architectural designs were explored, as elaborated in \autoref{section:ablations}. Notably, in cases where the model's receptive field is sufficiently restricted, the entropy model after training can be represented effectively with a compact lookup table.

\subsection{Entropy Threshold and Equalizing Context Length} For entropy-based patching models, we determine a patching threshold to produce a target average \textit{patch size} on the pretraining data mixture. In {\textsc BLT}, as opposed to token-based approaches, the \textit{patch size} is freely adjustable, which notably influences the model’s effective context size. To standardize the average context length and prevent an advantage from larger patch sizes, we fix the expected number of bytes per batch. Consequently, models configured with larger patch sizes operate with shorter sequences. For experiments with Llama 2 data, we set the context to 8k bytes. For models trained on the {\textsc BLT}-1T dataset, we use a mean context length of 16k bytes, while consistently maintaining a batch size of 16M bytes on average.

Although the batch size remains fixed, dynamic patching techniques lead to varying proportions of bytes per patch during data loading. To optimize efficiency, our implementation of {\textsc BLT} training organizes patches within batches such that padding in the computationally costly latent transformer is unnecessary. This guarantees a uniform patch count across all batches. During the training phase, byte sequences are either padded or, if needed, truncated to 12k and 24k bytes for the Llama 2 and {\textsc BLT}-1T datasets, respectively, mitigating memory surges caused by exceptionally large patch sequences.

\subsection{Entropy Model Context}
\label{section:entropy-context}
Our experiments demonstrate that entropy patching tends to produce increasingly larger patches in highly structured materials, such as those found in multiple choice assessments (refer to the MMLU example patching illustrated in \autoref{fig:mmlu_patching}), which characteristically exhibit substantial repetition. This behavior arises due to decreased entropy associated with repetitive elements present in the entropy model's context. Therefore, for the comprehensive deployment of {\textsc BLT}-Entropy with a patch size of 4.5, we refresh the entropy context by introducing new line characters and employ an approximate monotonicity constraint, as this approach mitigates the "entropy drift" linked to changes in context length. It should be noted that this modification pertains solely to entropy calculation methodology; the determination of the entropy threshold remains consistent with the established procedure.

\subsection{FLOPs Estimation} 
\label{section:flops}

Our calculation of transformer floating point operations is primarily based on the methodology outlined in Chinchilla~\citep{hoffmann2022training}, which accounts for flops from the feed-forward network, the qkvo projections in the self-attention module, and the costs associated with attention computation and the output projection. However, our approach diverges in that we model the input embedding layer as an efficient lookup operation rather than as a dense matrix multiplication, effectively making it a zero-flop process. Consistent with earlier studies, we assume that the backward pass requires double the flops of the forward pass.

In order to calculate flops \textit{per byte} for the {\textsc BLT} architecture, we aggregate the floating point operations required by the local encoder transformer, the global latent transformer, and the local decoder transformer, as well as those incurred by the cross attention modules present in both the encoder and decoder. The computation is formalized as follows:

\begin{align}
    \text{FL}_{\text{{\textsc BLT}}} &= \text{Transf. FL}(h_{\mathcal{G}}, l_{\mathcal{G}}, m=n_{ctx}/n_p, V=0) / n_p\\
   &+ \text{Transf. FL}(h_{\mathcal{E}}, l_{\mathcal{E}}, m=w_{\mathcal{E}}, V=0) \\
   &+ \text{Transf. FL}(h_{\mathcal{D}}, l_{\mathcal{D}}, m=w_{\mathcal{D}}, V=256) \\ 
    &+ \text{Cross Attn. FL}(h_{\mathcal{E}}, l_{\mathcal{E}}, m=n_p, r=n_p/k) \times k / n_p \\ 
   &+ \text{Cross Attn. FL}(h_{\mathcal{D}}, l_{\mathcal{D}}, m=k, r=k/n_p) 
\end{align}

Here, $n_{ctx}$ denotes the byte-level sequence length, while $n_p$ represents the size of each patch. The parameter $r$ specifies the proportion of queries relative to key/values, and $k$ defines the relationship between the patch and byte dimensions; specifically, it reflects how many local models are joined to constitute the overall model representation (with $k=2$ depicted in \autoref{fig:crossattn}). The symbol $V$ stands for the size of the output vocabulary, utilized solely within the local decoder. The context length $m$ employed in attention calculations varies depending on whether the operation is carried out on byte sequences or patch sequences. Consequently, we adjust the formulas for calculating attention FLOPs for each respective component. The precise formulations for both Transformer-FLOPs and Cross-Attention FLOPs are outlined in Appendix~\ref{section:flops-appendix}.

\subsection{Bits-Per-Byte Estimation} Perplexity is inherently tied to a particular tokenizer, since it quantifies prediction uncertainty at the level of individual tokens. To allow for meaningful comparisons between byte-level and token-level models, and in accordance with prior studies~\citep{xue2022byt5, yu2023megabyte, wang2024mambabyte}, we adopt the Bits-Per-Byte (BPB) metric, which serves as a tokenizer-agnostic counterpart to perplexity. Formally, BPB is computed as follows:

\begin{align}
    \text{BPB}(x)&= \frac{\mathcal{L}_{CE}(\pmb{x})}{\ln(2)\cdot n_{\text{bytes}}}
\end{align}

In this formulation, the overall uncertainty in $\pmb{x}$, given by the cumulative cross-entropy loss, is scaled by both the total count of bytes present in $\pmb{x}$ and a fixed normalization constant.

\subsection{Transformer Architecture Hyperparameters} In all transformer blocks within {\textsc BLT}, encompassing both local and global components, we adopt hyperparameter settings closely aligned with the Llama~3 architecture~\citep{dubey2024llama}. Specifically, our feed-forward networks employ the SwiGLU activation~\citep{swiglu}, and rotary positional embeddings (RoPE)~\citep{RoPe} with $\theta=500000$~\citep{xiong2024effective} are incorporated exclusively within self-attention modules. For normalization across layers, we apply RMSNorm~\citep{RMSNorm}. The implementation of all self-attention layers that utilize fixed-standard attention masks—such as \textit{block causal} and \textit{fixed-window block causal}—relies on Flash attention~\citep{dao2022flashattention}; when employing fixed-width attention masks, we set the window size to 512. Cross-attention layers, which necessitate dynamic, patch-dependent masking, are accelerated by Flex Attention\footnote{\url{https://pytorch.org/blog/flexattention}}, enabling fused operations and markedly enhancing training efficiency.

\subsection{{\textsc BLT}-Specific Hyperparameters} To evaluate the performance of {\textsc BLT} models, our experiments focus on two main aspects: scaling behavior and downstream task performance. Across both experimental directions, we examine models of varying sizes—specifically, 400M, 1B, 2B, 4B, and 8B parameters. Appendix~Table~\ref{tab:arch} provides the architectural hyperparameter settings for each scale.

Initialization of the queries for the initial cross-attention layer within the local encoder is performed using max-pooling. Across all {\textsc BLT} models, we utilize $500,000$ hashes with a single hash function, covering n-gram lengths between 3 and 8. Every model is trained using a learning rate of $4e-4$. This uniform learning rate selection is based on an extensive hyperparameter sweep in the $1e-3$ to $1e-4$ interval for 400M and 1B models, indicating that this rate is optimal for both token and {\textsc BLT} variants. For experiments concerning scaling on Llama-2 data, we adhere to batch sizes recommended by~\cite{dubey2024llama} or match the specified number of bytes. Training is carried out using the AdamW optimizer~\citep{adamw}, with parameters set as follows: $\beta_1=0.9$, $\beta_2=0.95$, and $\epsilon=10^{-8}$. We implement a learning rate schedule comprising 2000 steps of linear warm-up followed by cosine decay to zero, apply a weight decay factor of 0.1, and enforce global gradient clipping at a value of 1.0.

\section{Scaling Trends}
\label{section:scaling}

We offer a comprehensive analysis of the scaling behaviors of byte-level models, with the intention of guiding the further development and expansion of {\textsc BLT} models. Our study seeks to overcome shortcomings present in earlier works on byte-level models through several key approaches: (a) We examine scaling patterns within a compute-efficient training regime, (b) We develop and evaluate parallel 8B parameter models using substantial training datasets (reaching up to 1T tokens or 4T bytes) across various downstream benchmarks, and (c) We assess how scaling operates when inference costs are held constant. In subsequent sections, we will explore particular benefits unique to modeling at the byte-sequence level.

\subsection{Parameter Matched Compute Optimal Scaling Trends}
\label{section:parameter-matched-scaling}
We employ the Llama 2 dataset to train multiple \textit{compute-optimal} bpe and {\textsc BLT} models, spanning four model sizes from 1B to 8B parameters. Performance is assessed by plotting training flops against language modeling results on a representative subset of the training data mixture. For the bpe models, we utilize the optimal balance between model parameters and dataset size, following guidance from Llama~3~\citep{dubey2024llama}. This \textit{compute-optimal} configuration is postulated to maximize performance on the training data within a fixed computational budget~\citep{hoffmann2022training}, thereby serving as a strong reference point for our experiments. Correspondingly, for every bpe model, we construct a {\textsc BLT} model trained on identical data, with a Latent Transformer component matched in both scale and architectural details to its bpe Transformer counterpart.

As shown in~\autoref{fig:scaling-compute-opt} (right), {\textsc BLT} models consistently perform on par with or better than their bpe counterparts, and this pattern persists across varying model sizes and computational budgets. To our knowledge, {\textsc BLT} represents the first byte-level Transformer architecture to exhibit scaling behaviors equivalent to those of BPE-based models under compute-optimal conditions. This finding supports our hypothesis that the ideal parameter-to-compute ratio established for bpe is also applicable to {\textsc BLT}, or at least is quite similar.

Both enhancements in model architecture and the integration of dynamic patching are essential for aligning with bpe scaling patterns. In ~\autoref{fig:scaling-compute-opt} (left), we present a comparison between space-patching-derived models and Llama 3. SpaceByte \citep{slagle2024spacebyte} is estimated using {\textsc BLT} space-patching, excluding both n-gram embeddings and cross-attention mechanisms. While SpaceByte represents an advance relative to Megabyte, there remains a significant gap to Llama 3. In \autoref{fig:scaling-compute-opt} (right), we highlight the gains obtained through both modifications in model architecture and the adoption of dynamic patching. When evaluated at scale, {\textsc BLT} models achieve comparable performance to leading tokenizer-based architectures like Llama 3.

Additionally, we find that performance among tokenizer-based models is influenced by tokenizer selection: those trained with the Llama-3 tokenizer consistently surpass models utilizing the Llama-2 tokenizer when trained on identical datasets.

In summary, our {\textsc BLT} architecture exhibits scaling behavior that lies between Llama 2 and 3 when much larger patch sizes are utilized. While the bpe tokenizers for Llama 2 and 3 produce average token lengths of 3.7 and 4.4 bytes respectively, {\textsc BLT} can demonstrate comparable scaling characteristics with average patch sizes reaching 6 or even 8 bytes. Since inference flops are inversely related to the mean patch size, opting for a patch size of 8 bytes can yield close to a 50\% reduction in inference flops. Furthermore, models employing these larger patch sizes tend to see improved performance as both the model and dataset sizes increase. Notably, {\textsc BLT} configured with an 8-byte patch size begins with considerably lower performance relative to bpe Llama 2 at the 1B parameter scale, but surpasses bpe’s performance at the 7B scale. This pattern indicates that larger patch sizes might offer even greater advantages as models and training budgets are further scaled, and hints at the potential feasibility of still greater patch sizes under continued model growth.

\subsection{Beyond Compute Optimal Task Evaluations}
\label{section:evals}
To further investigate scaling behaviors, we continue training an 8B {\textsc BLT} model past the compute-optimal regime using the {\textsc BLT}-1T dataset, which is larger and contains higher-quality data. Performance is then assessed across a variety of established classification and generative benchmarks. For the evaluation process, we focus on tasks measuring common sense reasoning, factual knowledge, and code generation abilities:
\paragraph{Classification tasks} comprise ARC-Easy (0-shot)~\citep{Eval_ARC}, ARC-Challenge (0-shot)~\citep{Eval_ARC}, HellaSwag (0-shot)~\citep{Eval_hellaswag}, PIQA (0-shot)~\citep{Eval_piqa}, and MMLU (5-shot)~\citep{hendrycks2020measuring}. We use a prompt-scoring technique based on the probability assigned to each answer option and compute the mean accuracy for reporting.
\paragraph{Coding related generation tasks:}  We evaluate code synthesis by reporting pass@1 metrics on MBPP (3-shot)~\citep{Eval_MBPP} and HumanEval (0-shot)~\citep{Eval_HumanEval}, gauging how effectively the LLM can produce correct Python code solutions.

In \autoref{tab:evals}, we present a comparison of three models trained on the {\textsc BLT}-1T dataset: a bpe Llama 3 tokenizer-based model,\footnote{The Llama 3 tokenizer, with its 128k vocabulary, is selected as it yields better results compared to the 32k vocabulary of Llama 2.} alongside two variants of the {\textsc BLT} architecture. The first variant implements a space-patching method ({\textsc BLT}-Space), while the second applies an entropy-based patching approach ({\textsc BLT}-Entropy). The latter is configured with an approximate monotonicity constraint and its context is reset upon encountering new lines, as elaborated in~\autoref{section:entropy-context}. Each model is trained with a comparable computational (flop) budget. Notably, for {\textsc BLT}-Entropy, we also reduce the entropy threshold during inference from 0.6 to 0.1, observing that this modification enhances performance on tasks, albeit with a greater number of inference steps required.

The {\textsc BLT}-Entropy model surpasses the Llama 3 model in performance across 4 of the 7 evaluated tasks, despite both models being trained with an equivalent number of bytes. This advantage likely stems from two main factors: (1) more efficient utilization of training computation facilitated by dynamic patching, and (2) the explicit modeling of byte-level data rather than token-level representations.

Conversely, {\textsc BLT}-Space generally lags behind the Llama 3 tokenizer in all tasks except one; however, it offers notable improvements in inference efficiency due to its greater average patch size of 6 bytes. By contrast, both the bpe and entropy-patching approaches yield an average patch size close to 4.5 bytes when evaluated on the training data mixture. Under equivalent training resources, the model employing larger patch sizes is able to process 30\% more data than the bpe and entropy-patching models. This increased data coverage could, however, result in {\textsc BLT} deviating further from the point of optimal computational efficiency.

\subsection{Patches Scale Better Than Tokens} {\textsc BLT} models enable concurrent scaling of both model and patch sizes without altering the computational cost during training and inference, and without requiring more training data. Unlike token-based architectures, which are limited by a fixed vocabulary and associated efficiency constraints (see Section~\ref{section:bpe-patch}), patch-based models allow for arbitrarily large patches. This flexibility in increasing patch size reduces computation demands, thereby freeing up resources that can instead be utilized to expand the global latent transformer, since it needs to be executed less frequently.

We perform a fixed inference scaling analysis to examine whether increasing model size while reducing the number of steps on larger patches yields better results compared to smaller models taking more steps. Beginning with models containing 400 million and 3.6 billion parameters and utilizing the Llama 2 tokenizer, we identify flop-equivalent counterparts that use the Llama 3 tokenizer, as well as {\textsc BLT}-Entropy models with mean patch sizes of 6 and 8 bytes on the training datamix (refer to~\autoref{tab:fixed-inf-params} for further model specifications). For models with a patch size of 8, we employ 3 encoder layers as opposed to 1. All models are trained under multiple training flop constraints.

As illustrated in \autoref{fig:fixed_inference_scaling}, {\textsc BLT} models exhibit superior scaling behavior compared to tokenization-based models across both inference FLOP categories. Although BPE models initially outperform at lower training budgets, they are overtaken by {\textsc BLT} models shortly after exceeding the compute-optimal point. In real-world scenarios, allocating additional resources for the one-time pretraining phase can yield a model that performs better under a fixed inference budget. This approach is exemplified by the 8B parameter models, such as Llama 3.1, which have been exposed to training data volumes up to two orders of magnitude greater than the compute-optimal amount for that parameter size.

The point at which {\textsc BLT} surpasses token-based approaches has moved marginally closer to the compute-optimal regime with the transition to higher flop class models, shifting from three times to approximately 2.5 times the compute-optimal budget. Additionally, the model with patch size 8 exhibits a more pronounced scaling curve within the larger flop category, overtaking competing models at an earlier stage. As outlined in Section~\ref{section:parameter-matched-scaling}, models with increased patch sizes tend to approach the performance levels of BPE models as model scale grows. We partially ascribe this to the diminishing proportion of overall flops devoted to the byte-level Encoder and Decoder components, which do not scale as rapidly as the Latent Transformer. When expanding the total parameter count twentyfold from 400M to 8B, the local parameters within {\textsc BLT} are only roughly doubled. This observation is significant since enlarging patch sizes only influences the computational demands of the patch Latent Transformer, leaving the byte-level components unaffected. Indeed, this explains why the {\textsc BLT}-Entropy with ps=8 shifted from representing 1.6x to 1.7x the parameter count of the Llama 2 model upon scaling up to the larger model size.

In conclusion, our investigation into patch-length scaling indicates that the {\textsc BLT} patch-based architecture exhibits enhanced scaling performance when both the patch size and model size are increased in tandem. Notably, these positive scaling dynamics not only continue but may further strengthen as the model scale grows.

\section{Byte Modeling Improves Robustness}
We further evaluate the robustness of {\textsc BLT} relative to token-based models that do not incorporate explicit byte-level data, and introduce a method to convert pretrained token-based models into byte-based ones.
\label{sec:robustness}

\subsection{Character-Level Tasks}

An initial impetus for developing byte-level models stemmed from their inherent resilience to noise introduced at the byte level in inputs, as well as their sensitivity to the internal composition of tokens—a limitation often encountered in models that rely on tokenizers. To systematically assess these properties, we conduct supplementary experiments on benchmarks designed to test not only robustness to input noise, but also sensitivity to individual characters across both English and multiple languages, encompassing categories such as digits and phonemes. The outcomes of these evaluations are reported in Table~\ref{tab:char_tasks}.

\paragraph{Noisy Data} To assess the robustness of {\textsc BLT} relative to tokenizer-based models, we generate perturbed versions of the benchmark classification tasks outlined in Section~\ref{section:evals}. Specifically, we introduce variability using five different character-level noise mechanisms: (a) \textit{AntSpeak}: Transforms the entire input into uppercase letters with spaces between each character. (b) \textit{Drop}: Randomly eliminates 10\% of the characters within the text. (c) \textit{RandomCase}: For each character in the text, randomly assigns uppercase to 50\% and lowercase to the remaining 50\%. (d) \textit{Repeat}: Selects 20\% of characters and repeats each one up to four times. (e) \textit{UpperCase}: Converts all text characters to uppercase. When evaluating, we apply each noise type separately to the prompt, the completion, or both, and then compute the mean score across these settings. Results on the perturbed HellaSwag dataset~\citep{Eval_hellaswag}, presented in Table~\ref{tab:char_tasks}, reveal that {\textsc BLT} consistently surpasses tokenizer-based models in robustness, providing on average an 8-point improvement over models trained on the same data, and even outperforming Llama 3.1, which was trained on a much larger data corpus.

\paragraph{Phonology - Grapheme-to-Phoneme (G2P)} We evaluate how effectively {\textsc BLT} converts sequences of graphemes—that is, written characters that comprise a word—into corresponding phonemic representations indicating pronunciation. Table \ref{tab:char_tasks} displays the outcomes for the G2P task under a 5-shot evaluation setup, leveraging Phonology Bench~\citep{suvarna2024phonologybench}. Our findings demonstrate that {\textsc BLT} achieves superior performance compared to the baseline Llama 3 1T tokenizer-based model on this task.

\paragraph{CUTE}
To evaluate character-level capabilities, we test {\textsc BLT} on the CUTE benchmark~\citep{edman2024cute}, which includes a range of tasks generally grouped into three main areas: compositional understanding, recognizing orthographic similarity, and sequence manipulation skills. This suite of tasks presents a notable difficulty for models based on tokenization, as they often exhibit awareness of token spellings yet are less adept at leveraging this to effectively manipulate text at the character level. As summarized in Table~\ref{tab:char_tasks}, {\textsc BLT}-Entropy surpasses both BPE Llama 3 models by a margin exceeding 25 points. Notably, our model displays remarkable strength in character manipulation, attaining 99.9\% accuracy on both spelling-related tasks. This notable advancement is achieved even though {\textsc BLT} has been trained on merely one-sixteenth the data used for Llama 3.1, suggesting that BPE models inherently struggle with acquiring character-level competencies. Figure \ref{fig:cute_examples} presents several examples where the Llama 3 tokenizer model encounters difficulties but the {\textsc BLT} model excels. The tasks of word deletion and word insertion are the only ones in which BPE-based models have the upper hand. Such operations may pose challenges for byte-level models, yet the difference in performance is relatively small; moreover, constructing word-level representations from character-level information may be more tractable than attempting the reverse. For consistency, we follow an identical evaluation protocol across all tasks, employing the original Huggingface prompts. It is possible that BPE models could show enhanced outcomes with more refined prompt engineering.

\paragraph{Low Resource Machine Translation} We assess the capabilities of {\textsc BLT} for translation tasks involving both directions between English and twenty-one lower resource languages, representing six major language families and a variety of scripts, using the FLORES-101 benchmark~\citep{goyal2022flores}. The SentencePiece BLEU scores for these experiments are provided in Table~\ref{tab:flores}. The findings indicate that {\textsc BLT} consistently outperforms a comparable model utilizing the Llama 3 tokenizer, yielding an overall improvement of 2 BLEU points for translations into English, and a 0.5 point gain for translations from English. On widely-used language pairs, {\textsc BLT} achieves similar or marginally higher performance than Llama 3. Notably, for many language pairs within lower-resourced families, {\textsc BLT} surpasses Llama 3, highlighting the advantage of byte-level modeling in handling infrequent and diverse byte sequences.

\subsection{Training {\textsc BLT} from Llama 3}
\label{sec:distillation}
We investigate a methodology wherein {\textsc BLT} models benefit from prior work by utilizing pre-trained tokenizer-based models, facilitating improved and accelerated training convergence. This is implemented by initializing the global transformer parameters within {\textsc BLT} using those from a pre-trained Llama 3.1 model. The global transformer weights are subsequently fine-tuned using a learning rate set to one-tenth of that applied to the local encoder and decoder models, across the number of steps that is optimal for Llama 3. Our results, summarized in Table~\ref{tab:distill}, provide a comparison with a baseline {\textsc BLT}. The findings demonstrate that initializing {\textsc BLT} with Llama 3.1 parameters leads to marked improvements over both the Llama 3 and baseline {\textsc BLT} models, even though the total number of flops used in training is equivalent. Furthermore, relative to our {\textsc BLT}-Entropy variant (refer to Table~\ref{tab:evals}), which underwent training on a much larger corpus of 1T tokens, {\textsc BLT} initialized from Llama 3.1 continues to yield superior results on the MMLU task. This suggests that this approach can effectively cut down the computational cost associated with training.

This approach may be interpreted as transforming models reliant on tokenizers into ones that are tokenizer-independent, thereby adapting the pre-trained LLaMA 3.1 model into a {\textsc BLT} configuration. For thoroughness, Table \ref{tab:distill} presents a comparison between the original LLaMA 3.1—trained with 15T tokens—and its {\textsc BLT} counterpart. Our findings show that while there is only a minor decrease in performance on MMLU and HumanEval, other evaluation benchmarks display a more pronounced reduction. These results indicate the necessity for continued research to better exploit the capabilities of the pre-trained model, with particular emphasis on optimizing both the data composition and other relevant hyperparameters.

\section{Ablations and Discussion}
\label{section:ablations}
This section presents a series of ablation studies aimed at validating our design decisions for the {\textsc BLT} architecture, as well as the selection of patching methods and tuning of hyper-parameters used for the {\textsc BLT} model with 8B parameters, which is trained on the {\textsc BLT}-1T dataset.

\paragraph{Entropy Model Hyper-parameters} We examine the influence of entropy model capacity and the length of the context window on model scaling efficiency by training a set of byte-level entropy transformer models spanning model sizes from 1 million to 100 million parameters and using context windows between 64 and 512 tokens. In \autoref{fig:entropy_ablation}, we illustrate bpb as a function of training compute, reporting scaling law curves for the $400m$ and $1b$ parameter {\textsc BLT} models trained with the Llama-2 dataset. Results indicate that increasing either entropy model size or context window generally improves scaling, though the benefits taper off beyond 50 million parameters.

\paragraph{Types of Patching}

We conduct an ablation study on the four patching methods outlined in Section~\ref{section:patching}, specifically: 1) Strided Patching using strides of 4 and 6, 2) Whitespace-based Patching, 3) BPE Tokenizer patching aligned with the Llama 3 tokenizer, and 4) Entropy-driven patching utilizing a compact byte LLM.

Although dynamic patching shortens the actual sequence lengths, we regulate sequence length so that all patching strategies have a comparable context window. Each model processes an equivalent number of bytes per sequence during both training and inference, on average, thereby eliminating the possibility that improved performance is due to access to longer contexts. The results of these ablation studies are presented in Figure~\ref{fig:scaling-compute-opt}. Across these evaluations, every alternative patching scheme demonstrates superior performance compared to static patching, with space patching performing nearly as well as dynamic entropy-based patching.

In \autoref{tab:evals_ablation}, we report benchmark results for {\textsc BLT} models that utilize tokenizer-based methods, space patching, and entropy-based patching. These models were trained on the Llama 2 dataset for a tuned number of training steps~\citep{dubey2024llama}. While space patching employs a more straightforward approach, omitting the need to execute an entropy model dynamically during training, our findings indicate that the advantages demonstrated by entropy-based patching in scaling experiments~(Section~\ref{section:scaling}) extend to performance on downstream benchmark tasks as well.\footnote{Space patching outcomes are derived from initial runs lacking cross-attention, but comparable patterns are seen even when cross-attention is incorporated.}

\paragraph{Cross-Attention}
In~\autoref{tab:crossattn_ablations_1b}, we examine the effects of incorporating cross-attention at different layers within both the encoder and decoder components of {\textsc BLT}. Specifically, for encoder cross-attention, we explore three approaches to query initialization: 1) utilizing an identical learned embedding across all global states, 2) applying a hash embedding to the bytes comprising each patch, and 3) aggregating the encoder’s hidden states for the patch bytes at the relevant encoder layer.

Our results indicate that incorporating cross-attention within the \textit{decoder} yields the greatest benefits. Implementing cross-attention in the encoder offers modest gains, which are evident primarily when query pooling initialization is utilized. Furthermore, our analysis reveals that cross-attention proves advantageous on Common-Crawl datasets, with the improvements being most pronounced for larger patch sizes.

\paragraph{n-gram Hash Embeddings} 

We experiment with n-gram hash embedding vocabulary sizes of 0, 100K, 200K, and 400K, with the corresponding outcomes reported in Table~\ref{tab:ngram_ablations_1b}. The results demonstrate that the use of hash embeddings yields improvements across all evaluated domains, with the most significant effects observed on Wikipedia and Github, where a 0.04 bpb improvement is noted (in comparison to only 0.01 bpb after 15k steps at the 8B scale). When operating at the 8B model size, a reduction from 500K to 300K hashes results in a negligible performance drop of approximately 0.001 bpb after 15k steps. These findings suggest that hash embeddings are essential for enabling {\textsc BLT} to achieve performance parity with tokenizer-based approaches; nonetheless, benefits begin to plateau beyond 300K hashes. Moreover, the observed performance gains from hash embeddings are mostly additive to those provided by cross-attention mechanisms, as they offer enhancements on distinct datasets.

\paragraph{Local Model Hyperparameters} Table \ref{tab:local_layers} presents an ablation study examining different configurations for the number of layers in both the local encoder and decoder. When utilizing hash n-gram embeddings, {\textsc BLT} demonstrates strong performance when the encoder is highly minimal, such as employing a single layer, while the decoder is more substantial in complexity.

\section{Related Work}
\label{section:related_work}

\textbf{Character-Level RNNs:} Since the inception of neural networks, character-level language modeling has remained a widely investigated task~\citep{sutskever2011generating,mikolov2012subword,graves2013generating}, primarily due to its inherent capacity to seamlessly handle out-of-vocabulary terms, eliminating the need for explicit back-off mechanisms. \cite{kim2016character} proposed a model architecture that exclusively processes input characters by employing convolutional and highway layers, which subsequently provide inputs to LSTM-based RNNs; their approach achieved results on par with then state-of-the-art RNN language models for English, and surpassed them for morphologically complex languages—demonstrating an additional benefit of character-level LLMs. In the area of machine comprehension, \cite{kenter2018byte} utilized byte-level LSTM models that outperformed traditional word-level counterparts specifically for morphologically rich languages such as Turkish and Russian. Similarly, \cite{zhang2015character} explored classification tasks with character-based convolutional architectures, and found that they surpassed word-level models for select applications. Hierarchical LSTM models introduced by \citet{chung2019hierarchical}, which incorporate boundary detectors at each hierarchical layer, have also been shown to uncover underlying text structure and improve character-level language modeling performance. Meanwhile, \cite{kalchbrenner2016neural}'s ByteNet adopted character-level CNN-based layers in lieu of attention mechanisms within the context of machine translation.

\textbf{Character-Level Transformers:} The introduction of transformer architectures utilizing attention mechanisms~\citep{vaswani2017attention}, along with advances in subword tokenization~\citep{sennrich-etal-2016-neural}, led to substantial gains in the effectiveness of neural approaches for language modeling and standard evaluation tasks. Nevertheless, selecting word or subword units inherently establishes an inductive bias regarding the abstraction layer at which these models function. In an effort to merge the progress made by transformer-based models with the early encouraging outcomes of character-level language modeling, \citet{al2019character} employed extremely deep transformer networks. They incorporated auxiliary loss functions to successfully train transformer models that surpassed prior LSTM-based character language models. Despite these achievements, a notable disparity remained compared to word-level large language models. Similarly, GPT-2~\citep{radfordgpt2} reported that for extensive datasets such as the 1 billion word benchmark, language models operating at the byte level did not achieve results on par with their word-level counterparts.

Although \cite{Choe2019BridgingTG} showed that transformer-based byte-level LLMs can surpass subword-level LLMs with a similar number of parameters, these models require substantially more computational resources and training time. Likewise, \cite{el2020characterbert} introduced a BERT variant, CharFormer, which constructs word representations by applying convolutional layers to character embeddings and achieves performance gains in the medical domain, but at the cost of considerably greater computational expense. CANINE, proposed by \cite{clark2022canine}, is an encoder-only model with 150M parameters that processes raw character sequences directly. It incorporates a deep transformer stack akin to the core of our global model, combining a local transformer with strided convolutions for input character downsampling. This approach enables CANINE to surpass the performance of a similarly sized token-level encoder-only model (mBERT) across various downstream multilingual tasks. ByT5~\citep{xue2022byt5} investigates byte-level encoder-decoder architectures that omit patching mechanisms entirely. Their results indicate enhanced noise robustness and competitiveness with tokenizer-based models while requiring only a quarter of the training data; nevertheless, the absence of patching necessitates performing costly attention calculations on every byte, resulting in significant computational overhead. Handling input as raw bytes instead of subword units increases sequence lengths, thus posing additional challenges to scaling byte-level models efficiently. More recently, leveraging the Mamba Architecture~\citep{gu2023mamba}, which is capable of maintaining a constant-size memory state over long contexts, \cite{wang2024mambabyte} have trained a byte-level Mamba model—again eschewing patching—and achieved superior performance relative to byte-level transformers, as measured in bits-per-byte, on several datasets at the 350M parameter scale under flop-equivalent constraints.

\textbf{Patching-based approaches:} Leveraging patching methods has the potential to substantially reduce the number of flops required by byte-level LLMs while possibly maintaining model efficacy. Several studies have reported promising results at smaller scales regarding both model size and the quantity of training bytes. For example, \cite{nawrot-etal-2022-hierarchical} explore static patching strategies involving downsampling and upsampling, leading to the development of the hourglass transformer, which surpasses competing byte-level baselines at the 150M parameter level. Building on this, \cite{nawrot-etal-2023-efficient} introduce enhancements through dynamic patching techniques, such as a boundary predictor trained end-to-end, a boundary predictor guided by specific tokenizers, and an entropy-based patching mechanism resembling {\textsc BLT}. Their findings indicate that these methods outperform the standard transformers available at that time for language modeling at a scale of 40M parameters and approximately 400M training tokens. In a different vein, \cite{lester2024training} examine the impact of training on sequences that have been compressed via arithmetic coding, attaining greater compression rates than those achieved with BPE. By applying an equal-info windows method, they achieve superior results compared to byte-level baselines in language modeling; however, performance still falls short of subword-level baselines.

Our research is primarily influenced by and closely aligned with MegaByte~\citep{yu2023megabyte}, a causal, decoder-only LLM that relies on a fixed, static method for dividing input bytes into patches, concatenating their representations, and subsequently decoding patches back into bytes through a local model. The original MegaByte study established that this approach can achieve parity with tokenizer-based models when scaled to 1B parameters and evaluated on datasets of roughly 400B bytes. In our comprehensive ablations of MegaByte, we observe that its static patching scheme underperforms relative to the latest state-of-the-art, compute-efficient tokenizer-based models in flop-constrained scenarios, and we illustrate how {\textsc BLT} narrows this performance gap. Similarly, \cite{slagle2024spacebyte} identify this limitation in MegaByte, advocating for adaptive patching strategies—such as targeting whitespaces or analogous characters—and the addition of a local encoder module. Their findings highlight improvements over tokenized transformer models in compute-constrained evaluations within certain data domains, including Github and arXiv, at the 1B parameter level. We replicate experiments with this variant and show that further advances in model architecture are essential for byte-level approaches to attain true parity with the most advanced token-based models currently available, exemplified by systems like Llama 3.

\section{Limitations and Future Work}

For our architectural decisions, we employ models trained for the optimal number of steps as identified for Llama~3~\citep{dubey2024llama}. Nonetheless, it is important to note that these scaling laws are based on BPE-level transformers and may not yield the ideal ratios of data and parameter sizes when applied to {\textsc BLT}. Determining scaling laws tailored to {\textsc BLT}, which could reveal more advantageous scaling patterns for our approach, is reserved for future exploration. Furthermore, since many of our experiments utilize models with up to 1B parameters, it remains possible that the most effective architectural configurations may shift as model sizes reach 8B parameters or larger, potentially enhancing performance at such larger scales.

Current transformer codebases and libraries prioritize efficiency for architectures that utilize tokenizers. Although we report theoretical flop-equivalent experiments and adopt some efficient solutions—like FlexAttention—for layers differing from standard transformer designs, the real-world execution time of our implementations might still lag behind that of tokenizer-based models, indicating a potential for additional optimization.

Although {\textsc BLT} employs a separately trained entropy model for patching, integrating the patching model into an end-to-end training process presents a promising avenue for subsequent investigation. In Section \ref{sec:distillation}, we detail preliminary experiments that provide early evidence that "byte-ifying" models based on tokenizers—such as Llama 3, which has been trained on datasets exceeding 10T tokens—may be effective when the global transformer is initialized and its weights are frozen. Pursuing research in this direction could yield techniques that both maintain the advantages of bytefying and surpass the performance of existing tokenizer-based models, without necessitating training from the ground up.

\section{Conclusion}
\label{section:conclusion}

This work introduces the Byte Latent Transformer (\textbf{BLT}), an innovative model architecture that moves beyond the traditional reliance on fixed-vocabulary tokenization in large language models. {\textsc BLT} employs a flexible, learnable strategy for aggregating bytes into patches, thereby distributing computational effort according to the inherent complexity of the input data. This approach yields notable gains in both computational efficiency and model robustness. Through comprehensive scaling experiments, we show that {\textsc BLT} can achieve performance on par with established tokenization-based models such as Llama 3 when evaluated up to 8B parameters and 4T bytes. Moreover, it enables up to a 50\% reduction in inference computational requirements with only marginal drops in standard evaluation scores. Additionally, {\textsc BLT} introduces a novel avenue for scaling, supporting the joint expansion of both model capacity and patch size while maintaining a constant inference cost. This feature proves particularly beneficial under the computational constraints typical of real-world scenarios. By natively processing byte-level data, {\textsc BLT} enhances the model’s resilience to noisy or atypical inputs and provides a richer understanding of sub-word representations. Collectively, these advantages make {\textsc BLT} a strong candidate to supplant traditional tokenization-centric methods, offering a scalable, flexible, and robust foundation for future language models.

\end{document}